% GENERATED FILE. Edit canonical lessons, then run npm run generate:books.
% canonical-source-hash: fnv1a64:4705ad64ff7653c2
% canonical-lessons: RU-C349-sposobstvovat, RU-C349-privodit, RU-C349-vliyat, RU-C349-vliyanie, RU-C349-vozdeystvie

\chapter{Causes and influence}
\label{ch:ru-causes-and-influence}
\hlchaptermodality{349}

\begin{chapteropening}
\textbf{By the end of this chapter:} \emph{I can say that one thing helps bring about another, leads to it, influences it or has an effect on it.}

Complete the last lesson of chapter 349: I can say that one thing helps bring about another, leads to it, influences it or has an effect on it.
\end{chapteropening}

\section[\texorpdfstring{\ru{способствовать} (sposóbstvovat)}{способствовать (sposóbstvovat)}]{\ru{способствовать} (sposóbstvovat) — to contribute to, to promote}

\label{lesson:RU-C349-sposobstvovat}



\begin{quote}
Before the new one: say the Russian for a motive, then the Russian for origin.
\end{quote}

\subsection*{You'll want to know: \ru{способствовать}}
\textbf{\ru{способствовать}} (\emph{sposóbstvovat}) — \textquotedblleft{}to contribute to, to promote\textquotedblright{}.

\begin{quote}
\ru{Спорт} \ru{способствует} \ru{здоровому} \ru{сну}.
\end{quote}

\begin{quote}
\emph{Sport sposóbstvuyet zdoróvomu snu.}
\end{quote}

\begin{quote}
Sport promotes healthy sleep.
\end{quote}

Imperfective; present \emph{\ru{способствую}, \ru{способствуешь}}. It takes the dative: \emph{\ru{способствовать} \ru{успеху}}. Bookish.

\begin{cousinweb}
\textbf{\ru{способствовать}} is built on \emph{\ru{способ}} \textquotedblleft{}a way, a means\textquotedblright{} with the verb suffix \emph{-\ru{ствовать}}: to be a means to something. The further origin of \emph{\ru{способ}} is not settled.
\end{cousinweb}

\begin{grammarlens}[title={What you've built}]
One more word for telling, explaining and describing at B1.
\end{grammarlens}

\subsection*{Your turn}
\begin{itemize}
\raggedright
  \item \emph{Say it:} \emph{sposóbstvovat}
  \item \emph{Say it:} \emph{sposóbstvovat}, once more
  \item \emph{Say it:} a sentence of your own with \emph{sposóbstvovat}
  \item \emph{Recall:} say the Russian for a motive, then the Russian for origin, then say \emph{sposóbstvovat} again
\end{itemize}

\begin{tcolorbox}[breakable,colback=teal!4,colframe=teal!35!black,title={Before you move on}]
What does \ru{способствовать} mean? (\textquotedblleft{}To contribute to, to promote\textquotedblright{}.) Say it once more.
\end{tcolorbox}
\section[\texorpdfstring{\ru{приводить} (privodít)}{приводить (privodít)}]{\ru{приводить} (privodít) — to lead to, to result in}

\label{lesson:RU-C349-privodit}



\begin{quote}
Before the new one: say the Russian for origin, then the Russian for to contribute to.
\end{quote}

\subsection*{You'll want to know: \ru{приводить}}
\textbf{\ru{приводить}} (\emph{privodít}) — \textquotedblleft{}to lead to, to result in\textquotedblright{}.

\begin{quote}
\ru{Спешка} \ru{часто} \ru{приводит} \ru{к} \ru{ошибкам}.
\end{quote}

\begin{quote}
\emph{Spéshka chásto privódit k oshíbkam.}
\end{quote}

\begin{quote}
Haste often leads to mistakes.
\end{quote}

Imperfective; present \emph{\ru{привожу}, \ru{приводишь}}; perfective \emph{\ru{привести}}. \emph{\ru{Приводить} \ru{к}} + dative — to lead to. It also means to bring someone along: \emph{\ru{привести} \ru{друга}}.

\begin{cousinweb}
\textbf{\ru{приводить}} is built from \emph{\ru{при}-} \textquotedblleft{}to, towards\textquotedblright{} and \emph{\ru{водить}} \textquotedblleft{}to lead\textquotedblright{}, from Proto-Slavic \emph{*voditi} \textquotedblleft{}to lead, to conduct\textquotedblright{} (EtymDB). A cause leads to its result.
\end{cousinweb}

\begin{grammarlens}[title={What you've built}]
One more word for telling, explaining and describing at B1.
\end{grammarlens}

\subsection*{Your turn}
\begin{itemize}
\raggedright
  \item \emph{Say it:} \emph{privodít}
  \item \emph{Say it:} \emph{privodít}, once more
  \item \emph{Say it:} a sentence of your own with \emph{privodít}
  \item \emph{Recall:} say the Russian for origin, then the Russian for to contribute to, then say \emph{privodít} again
\end{itemize}

\begin{tcolorbox}[breakable,colback=teal!4,colframe=teal!35!black,title={Before you move on}]
What does \ru{приводить} mean? (\textquotedblleft{}To lead to, to result in\textquotedblright{}.) Say it once more.
\end{tcolorbox}
\section[\texorpdfstring{\ru{влиять} (vliyát)}{влиять (vliyát)}]{\ru{влиять} (vliyát) — to influence, to affect}

\label{lesson:RU-C349-vliyat}



\begin{quote}
Before the new one: say the Russian for to contribute to, then the Russian for to lead to.
\end{quote}

\subsection*{You'll want to know: \ru{влиять}}
\textbf{\ru{влиять}} (\emph{vliyát}) — \textquotedblleft{}to influence, to affect\textquotedblright{}.

\begin{quote}
\ru{Погода} \ru{влияет} \ru{на} \ru{моё} \ru{настроение}.
\end{quote}

\begin{quote}
\emph{Pogóda vliyáyet na moyó nastroyéniye.}
\end{quote}

\begin{quote}
The weather affects my mood.
\end{quote}

Imperfective; perfective \emph{\ru{повлиять}}. \emph{\ru{Влиять} \ru{на}} + accusative.

\begin{cousinweb}
\textbf{\ru{влиять}} was taken from Church Slavonic \emph{vŭlĭjati} (EtymDB) \textquotedblleft{}to pour in\textquotedblright{}, built from \emph{\ru{в}-} \textquotedblleft{}into\textquotedblright{} and the root of \emph{\ru{лить}} \textquotedblleft{}to pour\textquotedblright{}. To influence someone was to pour something into them.
\end{cousinweb}

\begin{grammarlens}[title={What you've built}]
One more word for telling, explaining and describing at B1.
\end{grammarlens}

\subsection*{Your turn}
\begin{itemize}
\raggedright
  \item \emph{Say it:} \emph{vliyát}
  \item \emph{Say it:} \emph{vliyát}, once more
  \item \emph{Say it:} a sentence of your own with \emph{vliyát}
  \item \emph{Recall:} say the Russian for to contribute to, then the Russian for to lead to, then say \emph{vliyát} again
\end{itemize}

\begin{tcolorbox}[breakable,colback=teal!4,colframe=teal!35!black,title={Before you move on}]
What does \ru{влиять} mean? (\textquotedblleft{}To influence, to affect\textquotedblright{}.) Say it once more.
\end{tcolorbox}
\section[\texorpdfstring{\ru{влияние} (vliyániye)}{влияние (vliyániye)}]{\ru{влияние} (vliyániye) — influence}

\label{lesson:RU-C349-vliyanie}



\begin{quote}
Before the new one: say the Russian for to lead to, then the Russian for to influence.
\end{quote}

\subsection*{You'll want to know: \ru{влияние}}
\textbf{\ru{влияние}} (\emph{vliyániye}) — \textquotedblleft{}influence\textquotedblright{}.

\begin{quote}
\ru{Родители} \ru{имеют} \ru{большое} \ru{влияние} \ru{на} \ru{детей}.
\end{quote}

\begin{quote}
\emph{Rodíteli iméyut bolshóye vliyániye na detéy.}
\end{quote}

\begin{quote}
Parents have a great influence on their children.
\end{quote}

Neuter. \emph{\ru{Под} \ru{влиянием}} + genitive — under the influence of.

\begin{cousinweb}
\textbf{\ru{влияние}} is built on \emph{\ru{влиять}} with the noun suffix \emph{-\ru{ние}}; the verb comes from Church Slavonic \emph{vŭlĭjati} \textquotedblleft{}to pour in\textquotedblright{} (EtymDB), on the root of \emph{\ru{лить}} \textquotedblleft{}to pour\textquotedblright{}.
\end{cousinweb}

\begin{grammarlens}[title={What you've built}]
One more word for telling, explaining and describing at B1.
\end{grammarlens}

\subsection*{Your turn}
\begin{itemize}
\raggedright
  \item \emph{Say it:} \emph{vliyániye}
  \item \emph{Say it:} \emph{vliyániye}, once more
  \item \emph{Say it:} a sentence of your own with \emph{vliyániye}
  \item \emph{Recall:} say the Russian for to lead to, then the Russian for to influence, then say \emph{vliyániye} again
\end{itemize}

\begin{tcolorbox}[breakable,colback=teal!4,colframe=teal!35!black,title={Before you move on}]
What does \ru{влияние} mean? (\textquotedblleft{}Influence\textquotedblright{}.) Say it once more.
\end{tcolorbox}
\section[\texorpdfstring{\ru{воздействие} (vozdéystviye)}{воздействие (vozdéystviye)}]{\ru{воздействие} (vozdéystviye) — an impact, an effect}

\label{lesson:RU-C349-vozdeystvie}



\begin{quote}
Before the new one: say the Russian for to influence, then the Russian for influence.
\end{quote}

\subsection*{You'll want to know: \ru{воздействие}}
\textbf{\ru{воздействие}} (\emph{vozdéystviye}) — \textquotedblleft{}an impact, an effect\textquotedblright{}.

\begin{quote}
\ru{Это} \ru{лекарство} \ru{оказывает} \ru{сильное} \ru{воздействие}.
\end{quote}

\begin{quote}
\emph{Éto lekárstvo okázyvayet sílnoye vozdéystviye.}
\end{quote}

\begin{quote}
This medicine has a strong effect.
\end{quote}

Neuter. \emph{\ru{Оказывать} \ru{воздействие} \ru{на}} + accusative — to have an effect on. Bookish.

\begin{cousinweb}
\textbf{\ru{воздействие}} is built from \emph{\ru{воз}-} and \emph{\ru{действие}} \textquotedblleft{}action\textquotedblright{}, from \emph{\ru{действовать}} \textquotedblleft{}to act\textquotedblright{}; the root is usually traced to the old verb \emph{\ru{деять}} \textquotedblleft{}to do\textquotedblright{}.
\end{cousinweb}

\begin{grammarlens}[title={What you've built}]
One more word for telling, explaining and describing at B1.
\end{grammarlens}

\subsection*{Your turn}
\begin{itemize}
\raggedright
  \item \emph{Say it:} \emph{vozdéystviye}
  \item \emph{Say it:} \emph{vozdéystviye}, once more
  \item \emph{Say it:} a sentence of your own with \emph{vozdéystviye}
  \item \emph{Recall:} say the Russian for to influence, then the Russian for influence, then say \emph{vozdéystviye} again
\end{itemize}

\begin{tcolorbox}[breakable,colback=teal!4,colframe=teal!35!black,title={Before you move on}]
What does \ru{воздействие} mean? (\textquotedblleft{}An impact, an effect\textquotedblright{}.) Say it once more.
\end{tcolorbox}
